\documentclass[12pt,a4paper]{article}
\usepackage[margin=1in]{geometry}
\usepackage{amsmath,amssymb,amsthm}
\usepackage{enumitem}
\usepackage{booktabs}
\usepackage{array}

\newtheorem{theorem}{Theorem}[section]
\newtheorem{corollary}[theorem]{Corollary}
\newtheorem{lemma}[theorem]{Lemma}
\newtheorem{definition}[theorem]{Definition}
\theoremstyle{remark}
\newtheorem*{remark}{Remark}

\title{Part 0 --- Supplement\\[6pt]
\large Inter-Layer Coupling and the Independent-Step Correction}
\author{RTAC}
\date{March 2026}

\begin{document}
\setlength{\emergencystretch}{3em}
\maketitle

\noindent This supplement extends Paper 0 by computing the inter-layer correction to the cascade's independent-step approximation. The slicing recurrence treats each compactification step as independent: the propagator through $n$ steps is the product of $n$ individual lapses. This is exact for volumes (the telescoping identity $V_D/V_4 = \prod N(d)$ holds by the Gamma function). It is not exact for quantities that depend on the overlap structure of consecutive layer integrands. The correction is determined by the eigenvalue deficit of the Gram correlation matrix of the integrands $f_d(x) = (1 - x^2)^{d/2}$---a Beta function identity with no free parameters.

The Gram correlation has an exact closed form in cascade-native primitives (Theorem~\ref{thm:gram-R}): $C^2_{d,d+1} = R(2d+2)^2 / [R(2d+1)\,R(2d+3)]$, equivalently $\log C^2_{d,d+1} = -\tfrac{1}{2}\Delta^2 \log\alpha\big|_{2d+2}$ (Corollary~\ref{cor:gram-laplacian}). The Gram correction is therefore the discrete log-Laplacian of the cascade scalar action's compliance function $\alpha(d) = R(d)^2/4$, evaluated at the doubled cascade argument. This connects the supplement directly to the action principle of Paper~IVb Remark~4.6: the supplement's path-distributed correction and Paper~IVb's single-source $\alpha(d^*)/\chi^k$ family are complementary derivatives of the same compliance function---linear source response (Paper~IVb) and discrete log-Laplacian (this supplement).

Throughout: ``second order'' refers exclusively to the perturbation-theory order in the eigenvalue expansion (Theorem~\ref{thm:pert}).

\setcounter{section}{14}
\section{Inter-Layer Coupling}

\subsection{The Gram matrix of cascade layer integrands}

\begin{definition}[Gram matrix]
For a set of cascade layers $\{d_1, d_2, \ldots, d_n\}$, the Gram matrix is
\[
G_{ij} = \langle f_{d_i}, f_{d_j} \rangle = \int_{-1}^{1} (1 - x^2)^{(d_i + d_j)/2}\, dx.
\]
\end{definition}

\begin{theorem}[Gram matrix as Beta function]\label{thm:gram}
\[
G_{ij} = B\!\left(\tfrac{1}{2},\, \tfrac{d_i + d_j}{2} + 1\right) = \sqrt{\pi}\, \frac{\Gamma\!\left(\frac{d_i + d_j}{2} + 1\right)}{\Gamma\!\left(\frac{d_i + d_j}{2} + \frac{3}{2}\right)}.
\]
\end{theorem}

\begin{proof}
Substituting $t = x^2$ in the integral: $\int_{-1}^{1}(1 - x^2)^\alpha\, dx = B(1/2, \alpha + 1)$ with $\alpha = (d_i + d_j)/2$. The Beta function identity $B(a,b) = \Gamma(a)\Gamma(b)/\Gamma(a+b)$ with $a = 1/2$ gives the stated form.
\end{proof}

\begin{definition}[Correlation matrix]
The normalised Gram matrix (correlation matrix) is
\[
C_{ij} = \frac{G_{ij}}{\sqrt{G_{ii}\, G_{jj}}} = \frac{B\!\left(\frac{1}{2},\, \frac{d_i + d_j}{2} + 1\right)}{\sqrt{B\!\left(\frac{1}{2},\, d_i + 1\right)\, B\!\left(\frac{1}{2},\, d_j + 1\right)}}.
\]
$C_{ii} = 1$ for all $i$, $\mathrm{tr}\, C = n$, and $0 < C_{ij} < 1$ for $i \neq j$.
\end{definition}

\subsection{The adjacent-layer overlap deficit}

\begin{theorem}[Overlap deficit]\label{thm:overlap}
For adjacent layers $d$ and $d+1$:
\[
1 - C_{d,d+1}^2 = 1 - \frac{B\!\left(\frac{1}{2},\, d + \frac{3}{2}\right)^2}{B\!\left(\frac{1}{2},\, d+1\right)\, B\!\left(\frac{1}{2},\, d+2\right)}.
\]
This quantity is strictly positive for all $d \geq 1$ and satisfies:
\[
1 - C_{d,d+1}^2 \sim \frac{1}{8\, d^2} \quad (d \to \infty).
\]
Numerically: $1 - C_{5,6}^2 = 0.00319$, $\; 1 - C_{11,12}^2 = 0.00083$, $\; 1 - C_{19,20}^2 = 0.00030$.
\end{theorem}

\begin{proof}
\emph{Positivity.} $C_{d,d+1} < 1$ by Cauchy--Schwarz: the integrands $f_d$ and $f_{d+1}$ are not proportional in $L^2[-1,1]$ for $d \neq d+1$, so the normalised inner product is strictly less than~$1$.

\emph{Asymptotic $1 - C_{d,d+1}^2 \sim 1/(8d^2)$.} Set $h(\alpha) = B(1/2,\alpha+1) = \sqrt{\pi}\,\Gamma(\alpha+1)/\Gamma(\alpha+3/2)$ and $L(\alpha) = \log h(\alpha)$. Then
\[
\log C_{d,d+1}^2 \;=\; 2L(d+\tfrac{1}{2}) - L(d) - L(d+1),
\]
the centered second difference of $L$ at $d+1/2$ with step $1/2$. Taylor expanding,
\[
L(d) + L(d+1) \;=\; 2L(d+\tfrac{1}{2}) + \tfrac{1}{4}L''(d+\tfrac{1}{2}) + O(L^{(4)}/192),
\]
so $\log C_{d,d+1}^2 = -\tfrac{1}{4}L''(d+\tfrac{1}{2}) + O(L^{(4)})$.

Now $L'(\alpha) = \psi(\alpha+1) - \psi(\alpha+3/2)$, where $\psi$ is the digamma function. The asymptotic expansion $\psi(x) = \log x - 1/(2x) - 1/(12x^2) + O(1/x^4)$ gives, after expanding $\psi(\alpha+1)$ and $\psi(\alpha+3/2)$ around $\alpha$ and combining,
\[
L'(\alpha) \;=\; -\frac{1}{2\alpha} + \frac{3}{8\alpha^2} + O(1/\alpha^3),
\qquad
L''(\alpha) \;=\; \frac{1}{2\alpha^2} + O(1/\alpha^3).
\]
Therefore
\[
\log C_{d,d+1}^2 \;=\; -\frac{1}{4}\cdot\frac{1}{2d^2} + O(1/d^3) \;=\; -\frac{1}{8d^2} + O(1/d^3),
\]
and $1 - C_{d,d+1}^2 = -\log C_{d,d+1}^2 + O\bigl((\log C^2)^2\bigr) = 1/(8d^2) + O(1/d^3)$.

The numerical values quoted in the theorem statement ($1 - C_{5,6}^2 = 0.00319$, $1 - C_{11,12}^2 = 0.00083$, $1 - C_{19,20}^2 = 0.00030$) are evaluated directly from the Beta function ratio; the asymptotic $1/(8d^2)$ supplies the leading $1/d^2$ scaling and the sign, but subleading $1/d^3$ terms matter at small $d$ (the asymptotic over-estimates by $\sim 50\%$ at $d=5$, by $\sim 25\%$ at $d=11$, and by $\sim 15\%$ at $d=19$, converging to the exact value as $d\to\infty$).
\end{proof}

\begin{remark}
The overlap deficit $1 - C_{d,d+1}^2$ is the geometric content at step $d$ that is \emph{not} shared with step $d+1$. It measures the imperfect folding: the fraction by which consecutive layer integrands fail to be collinear in $L^2$. Its positivity is a consequence of $R_{\mathrm{eff}}(d) = 1/\sqrt{d+3} > 0$ (Theorem~4.2): the compactification at each step is asymptotic, never complete.
\end{remark}

\subsection{Closed form via cascade slicing ratios}\label{sec:gram-closed-form}

The asymptotic of Theorem~\ref{thm:overlap} extends to an \emph{exact} closed form expressing $C^2_{d,d+1}$ in cascade-native primitives. The exact form reveals that the Gram correlation is the discrete log-Laplacian of the cascade scalar action's compliance function $\alpha(d) = R(d)^2/4$, evaluated at the doubled cascade argument $2d+2$.

\begin{theorem}[Gram correlation in cascade slicing ratios]\label{thm:gram-R}
For all $d \geq 1$, the Gram correlation between adjacent cascade layers is
\begin{equation}\label{eq:gram-R-form}
C^2_{d,d+1} \;=\; \frac{R(2d+2)^2}{R(2d+1)\,R(2d+3)},
\end{equation}
where $R(d) = \Gamma((d+1)/2)/\Gamma((d+2)/2)$ is the cascade slicing ratio of Paper~0 \S2.
\end{theorem}

\begin{proof}
The Gram entries (Theorem~\ref{thm:gram}) are Beta functions $G_{ij} = B(1/2, (d_i+d_j)/2 + 1)$. Using the Beta--Gamma reduction $B(1/2, n+1) = \sqrt{\pi}\,\Gamma(n+1)/\Gamma(n+3/2)$ together with the cascade slicing ratio's defining identity $R(d) = \Gamma((d+1)/2)/\Gamma((d+2)/2)$, one obtains:
\[
G_{d,d}     = B(\tfrac{1}{2}, d+1)   = \sqrt{\pi}\,R(2d+1), \qquad
G_{d,d+1}   = B(\tfrac{1}{2}, d+\tfrac{3}{2}) = \sqrt{\pi}\,R(2d+2),
\]
\[
G_{d+1,d+1} = B(\tfrac{1}{2}, d+2)   = \sqrt{\pi}\,R(2d+3).
\]
Substituting into the definition $C^2_{d,d+1} = G_{d,d+1}^2 / (G_{d,d}\,G_{d+1,d+1})$ gives
$C^2_{d,d+1} = R(2d+2)^2 / [R(2d+1)\,R(2d+3)]$.
\end{proof}

\begin{corollary}[Gram-Laplacian identity]\label{cor:gram-laplacian}
The log-correlation between adjacent cascade layers is the discrete log-Laplacian of the cascade scalar action's compliance function $\alpha(d) = R(d)^2/4$, evaluated at the doubled argument $2d+2$ with unit step:
\begin{equation}\label{eq:gram-laplacian}
\log C^2_{d,d+1} \;=\; -\frac{1}{2}\,\Delta^2 \log\alpha\,\Big|_{2d+2},
\end{equation}
where $\Delta^2 f|_n := f(n-1) + f(n+1) - 2f(n)$ is the centered discrete Laplacian.
\end{corollary}

\begin{proof}
By Theorem~\ref{thm:gram-R} and properties of the logarithm,
\[
\log C^2_{d,d+1} = 2\log R(2d+2) - \log R(2d+1) - \log R(2d+3) = -\Delta^2 \log R\,\big|_{2d+2}.
\]
Since $\alpha(d) = R(d)^2/4$, we have $\log\alpha = 2\log R - 2\log 2$, so $\Delta^2 \log\alpha = 2\,\Delta^2 \log R$ (constants drop in second differences). Substituting yields equation~\eqref{eq:gram-laplacian}.
\end{proof}

\begin{remark}[Numerical verification]\label{rem:gram-verifier}
Both Theorem~\ref{thm:gram-R} (the closed form) and Corollary~\ref{cor:gram-laplacian} (the Laplacian identity) are verified at all tested $d$ values in
\begin{center}
\texttt{tools/verifiers/gram\_compliance\_laplacian.py}
\end{center}
to machine precision on the closed form and to $\sim 10^{-7}$ relative on the Laplacian identity (the floating-point cancellation noise floor at large $d$, where $\log C^2_{d,d+1} \sim 1/(8d^2)$ is computed as a difference of $O(1)$ log-Gamma terms).
\end{remark}

\begin{remark}[Structural unification with the cascade scalar action]\label{rem:gram-action-unification}
Corollary~\ref{cor:gram-laplacian} establishes a direct structural connection between the supplement's Gram framework and the cascade scalar action $S[\varphi] = \sum_d (2\alpha(d))^{-1}(\Delta\varphi)^2$ of Paper~IVb Remark~4.6. Both correction families derive from the same compliance function $\alpha(d) = R(d)^2/4$:
\begin{itemize}
\item Paper~IVb's $\alpha(d^*)/\chi^k$ family: \emph{linear} source response. The cascade action's marginal Green's function at distinguished source layer $d^*$ equals $\alpha(d^*)$ exactly (Paper~IVb's marginal Green's function identity); chirality factorisation through $k$ independent channels gives the $1/\chi^k$ attenuation.
\item This supplement's adjacent overlap deficit: \emph{discrete log-Laplacian} of the same compliance function, at the doubled cascade argument. The deficit per step is $\frac{1}{2}\Delta^2(-\log\alpha)|_{2d+2}$ (positive because $-\log\alpha$ is convex on the cascade tower, i.e.\ $\alpha$ is log-concave).
\end{itemize}
The two correction families are complementary in the action's dual variables: source response (zeroth-order, single-layer) versus log-compliance curvature (second-order, distributed across the path). The doubling $d\mapsto 2d+2$ is a Beta--Gamma reduction artefact: cascade integrand inner products $\langle f_d, f_{d'}\rangle$ produce $R$-values at doubled arguments via $B(1/2, n+1) = \sqrt{\pi}\,R(2n+1)$ and $B(1/2, n+3/2) = \sqrt{\pi}\,R(2n+2)$, so the Gram correlation between layers $d$ and $d+1$ depends on cascade $R$-values at doubled indices $\{2d+1, 2d+2, 2d+3\}$.

This corollary closes Paper~I's open question on first-principles derivation of the Gram correction (Paper~I, ``Does not'' list, item~3): the derivation is the cascade compliance Laplacian.
\end{remark}

\subsection{The eigenvalue structure of the correlation matrix}

\begin{theorem}[Near-collinearity]\label{thm:eigenvalue}
For a path of $n$ consecutive layers $\{d_0, d_0+1, \ldots, d_0+n-1\}$, the correlation matrix $C$ has one dominant eigenvalue $\lambda_1$ satisfying
\[
\frac{\lambda_1}{n} = 1 - \epsilon(n, d_0), \qquad \epsilon > 0,
\]
and $(n - 1)$ subdominant eigenvalues summing to $n\,\epsilon$. The deficit $\epsilon$ depends on both $n$ (the number of layers) and $d_0$ (where in the cascade the path begins).
\end{theorem}

\begin{proof}
The entries $C_{ij}$ are smooth functions of $|d_i - d_j|$ that decrease from $C_{ii} = 1$ towards a minimum at $|d_i - d_j| = n - 1$. For the path $d = 5, \ldots, 12$ ($n = 8$): the minimum entry is $C_{5,12} = 0.9622$, far from zero. The matrix is therefore close to the rank-1 matrix $\mathbf{1}\mathbf{1}^T/n$, whose sole nonzero eigenvalue is $n$. The deficit $\epsilon$ is determined by the off-diagonal decay rate of $C_{ij}$, which is controlled by the Beta function asymptotics of Theorem~\ref{thm:overlap}.
\end{proof}

\noindent Numerical eigenvalue structure for $d = 5, \ldots, 12$ ($n = 8$):

\begin{center}
\renewcommand{\arraystretch}{1.2}
\begin{tabular}{ccc}
\toprule
Eigenvalue & Value & Fraction of trace \\
\midrule
$\lambda_1$ & 7.9347263 & 99.184\% \\
$\lambda_2$ & 0.0646528 & 0.808\% \\
$\lambda_3$ & 0.0006152 & 0.008\% \\
$\lambda_4$--$\lambda_8$ & $< 10^{-4}$ & $< 0.001\%$ \\
\midrule
$\epsilon = 1 - \lambda_1/n$ & \multicolumn{2}{c}{$0.008159 = 0.816\%$} \\
\bottomrule
\end{tabular}
\end{center}

\noindent The cascade layers are 99.2\% collinear. The remaining 0.8\% is the inter-layer coupling.

\subsection{Path dependence of the eigenvalue deficit}

The deficit $\epsilon(n, d_0)$ depends on where in the cascade the path lies. Lower $d$ (closer to the observer, larger $R_{\mathrm{eff}}$, less complete folding) gives larger deficits.

\begin{center}
\renewcommand{\arraystretch}{1.2}
\begin{tabular}{llcc}
\toprule
Path & $n$ & $\epsilon = 1 - \lambda_1/n$ & Needed correction \\
\midrule
$d = 5, \ldots, 12$ & 8 & 0.00816 & $1.73\%$ ($\alpha_s$) \\
$d = 6, \ldots, 13$ & 8 & 0.00657 & $1.75\%$ ($m_\tau/m_\mu$) \\
$d = 10, \ldots, 17$ & 8 & 0.00332 & --- \\
$d = 14, \ldots, 21$ & 8 & 0.00200 & $0.13\%$ ($m_\mu/m_e$) \\
\bottomrule
\end{tabular}
\end{center}

\noindent At fixed $n$, $\epsilon$ decreases with $d_0$ because the overlap deficit per step (Theorem~\ref{thm:overlap}) decreases as $1/d^2$.

\subsection{The independent-step correction}

The cascade's leading-order predictions treat each compactification step as independent. A quantity that depends on the cascade's geometry over a path of $n$ layers receives a correction from the inter-layer coupling. Two distinct measures of this coupling are derived from the Gram matrix:

\begin{itemize}
\item \emph{Eigenvalue deficit $\epsilon = 1 - \lambda_1/n$} (Theorem~\ref{thm:eigenvalue}): the matrix-level magnitude of inter-layer coupling, characterised by the dominant eigenmode of $C$ and aggregated over all $n^2$ pair correlations with weight $(1-C_{ij})$. Derived perturbatively in Section~\ref{sec:perturbation}.
\item \emph{Geometric correction $\delta\Phi$} (Definition~\ref{def:delta-phi} below): the multiplicative log-decrement a propagator picks up traversing the path, summed over adjacent-layer overlap deficits with weight $(1-C^2_{d,d+1})$. Derived in Section~\ref{sec:first-order}.
\end{itemize}

The two are related but structurally distinct: $\epsilon$ characterises the Gram matrix as a whole; $\delta\Phi$ is the cumulative shift along a specific path. The corrected observable formula uses $\delta\Phi$.

\begin{definition}[Geometric correction along a path]\label{def:delta-phi}
For a cascade observable traversing the path $\{d_0, d_0+1, \ldots, d_0+n-1\}$,
\[
\delta\Phi \;:=\; \sum_{k=0}^{n-2}\bigl(1 - C^2_{d_k, d_{k+1}}\bigr),
\]
the path-cumulative adjacent overlap deficit. By Theorem~\ref{thm:overlap}, every term is strictly positive, so $\delta\Phi > 0$ for every cascade path. Theorem~\ref{thm:first-order} below shows that $\delta\Phi$ equals the first-order multiplicative correction to a cascade-descent observable, and Corollary~\ref{cor:exp-resum} shows that the unique closed-form correction respecting path-splitting composition is $Q = Q_0 \cdot \exp(\delta\Phi)$.
\end{definition}

\begin{remark}[What the correction is not]
The correction is not to the Gamma function, which is exact. It is not to the slicing recurrence, which is exact (the telescoping identity $V_D/V_4 = \prod N(d)$ holds by $\Gamma(a+1) = a\Gamma(a)$). It is to the \emph{independent-step approximation}: the assumption that a quantity depending on $n$ cascade layers is the product of $n$ individually exact factors. The layers are individually exact but not independent---their integrands overlap in $L^2$, and the eigenvalue deficit $\epsilon$ quantifies the overlap.
\end{remark}

\subsection{The perturbation theory for $\epsilon$}\label{sec:perturbation}

\begin{theorem}[First-order eigenvalue deficit]\label{thm:pert}
To first order in the deficit matrix $D = J - C$,
\[
\epsilon \;=\; 1 - \frac{\lambda_1}{n} \;=\; \frac{2}{n^2} \sum_{i < j} \left(1 - C_{ij}\right) \;+\; O\!\bigl(\|D\|^2\bigr),
\]
where $C_{ij}$ are the entries of the correlation matrix. The second-order correction is \emph{positive} (a sum of squared matrix elements of $D$ in the zero-eigenspace of $J$; see proof), so the right-hand side is an \emph{upper bound} for the eigenvalue deficit $\epsilon$.
\end{theorem}

\begin{proof}
Write $C = J - D$ where $J_{ij} = 1$ for all $i,j$ and $D_{ij} = 1 - C_{ij} \geq 0$, with $D_{ii} = 0$. The matrix $J$ has the non-degenerate eigenvalue $n$ with eigenvector $u = (1, \ldots, 1)/\sqrt{n}$, and eigenvalue $0$ with multiplicity $n - 1$ (eigenspace $u^\perp$). By standard Rayleigh--Schr\"odinger perturbation theory for the non-degenerate eigenvalue $n$:
\[
\lambda_1(C) \;=\; n - \langle u | D | u \rangle \;+\; \frac{1}{n}\sum_{v \in u^\perp}|\langle v | D | u \rangle|^2 \;+\; O\!\bigl(\|D\|^3\bigr).
\]
The first-order term evaluates to
\[
\langle u | D | u \rangle \;=\; \frac{1}{n}\sum_{i,j}D_{ij} \;=\; \frac{2}{n}\sum_{i<j}(1 - C_{ij}),
\]
giving the stated first-order formula after dividing by $n$. The second-order term is a sum of squared matrix elements, hence non-negative; it vanishes only when $D$ has no component in the $u^\perp$ sector (the rank-1 case, when $C$ is a constant-off-diagonal matrix). In all non-trivial cases it is strictly positive, so $\lambda_1(C) > n - \langle u|D|u\rangle$ at second order, equivalently $\epsilon < (2/n^2)\sum_{i<j}(1 - C_{ij})$: the first-order formula over-estimates the deficit, and the right-hand side is an upper bound.
\end{proof}

\noindent Verification: for $d = 5, \ldots, 12$ ($n = 8$), the first-order formula gives $\epsilon^{(1)} = 0.008173$; the exact eigenvalue gives $\epsilon = 0.008159$. Agreement: $0.17\%$; the exact value is below the first-order formula as predicted by the upper-bound direction of the theorem. The formula is verified to comparable precision for all paths tested.

\begin{corollary}[Closed form for $\lambda_1$ at $n = 2$]\label{cor:lambda1-n2}
For a two-layer path $\{d_0, d_0 + 1\}$, the dominant eigenvalue admits the exact closed form
\[
\lambda_1(2, d_0) \;=\; 1 + \frac{R(2 d_0 + 2)}{\sqrt{R(2 d_0 + 1)\,R(2 d_0 + 3)}}
\;=\; 1 + \exp\!\Bigl(-\tfrac{1}{4}\,\Delta^2 \log\alpha\big|_{2d_0+2}\Bigr),
\]
where the second equality follows from Corollary~\ref{cor:gram-laplacian}. Equivalently, $\epsilon(2, d_0) = 1 - \lambda_1/2 = (1 - C_{d_0, d_0+1})/2$, and the perturbation-theory upper bound of Theorem~\ref{thm:pert} is saturated at $n = 2$.
\end{corollary}

\begin{proof}
At $n = 2$ the matrix $C = \begin{pmatrix}1 & c \\ c & 1\end{pmatrix}$ with $c = C_{d_0, d_0 + 1}$ has eigenvectors $u_\pm = (1, \pm 1)/\sqrt{2}$ and eigenvalues $1 \pm c$, so $\lambda_1 = 1 + c$. The uniform vector $u = u_+$ is exactly the dominant eigenvector; the second-order correction in the proof of Theorem~\ref{thm:pert} vanishes because $u^\perp$ is one-dimensional and orthogonal to $u$ by the symmetry of $C$. Substituting Theorem~\ref{thm:gram-R} (and its Laplacian form, Corollary~\ref{cor:gram-laplacian}) for $c$ gives the stated closed form.
\end{proof}

\begin{remark}[Status of the $\lambda_1$ closed form for $n \geq 3$]\label{rem:lambda1-no-closed-form}
For $n \geq 3$ no single-equation radical form for $\lambda_1$ in cascade slicing ratios has been identified. What is available is the power-trace bound family $U_k = (\mathrm{tr}\,C^k)^{1/k}$ (Corollary~\ref{cor:lambda1-power-trace}), which provides closed-form approximations to $\lambda_1$ converging at geometric rate; combined with the $n = 2$ exact identity (Corollary~\ref{cor:lambda1-n2}) and the perturbation-theory bound of Theorem~\ref{thm:pert}, these cover all cascade-physics paths to better than machine precision.
\end{remark}

\begin{remark}[OP factorisation of the cascade correlation matrix]\label{rem:lambda1-OP-factorisation}
$C$ factorises as $C = A A^{\top}$ with $A$ lower triangular and $A_{ik} = \langle p_k, x^i\rangle_{\mu_m} / \|x^i\|_{\mu_m}$, where $\{p_k\}$ are the orthonormal polynomials of the measure $d\mu_m(x) = x^m/\sqrt{1-x^2}\,dx$ on $[0,1]$ ($m = 2 d_0 + 1$). Hence $\lambda_k(C) = \sigma_k(A)^2$. Under $u = x^2$, $\mu_m$ is a shifted Jacobi measure, identifying the closed-form question for $\lambda_1$ as a question in classical orthogonal-polynomial theory (\texttt{tools/research/lambda1\_jacobi\_route.py}).
\end{remark}

\begin{corollary}[Power-trace bound family closing $\lambda_1$ to arbitrary precision]\label{cor:lambda1-power-trace}
For each $k \geq 1$,
\[
U_k \;=\; \bigl(\mathrm{tr}\,C^k\bigr)^{1/k} \;\geq\; \lambda_1,
\]
with $U_1 \geq U_2 \geq \cdots \geq \lambda_1$ monotone, and $U_k \to \lambda_1$ at geometric rate $(\lambda_2/\lambda_1)^k$. Each $U_k$ is closed-form in cascade slicing ratios:
\[
\mathrm{tr}\,C^k \;=\; \sum_{i_1, \ldots, i_k} C_{i_1 i_2}\,C_{i_2 i_3} \cdots C_{i_k i_1},
\]
a polynomial sum of $R$-ratio products. Numerically, for cascade-physics paths ($d_0 \in \{5, \ldots, 19\}$, $n \leq 10$): $U_2$ is tight to $\sim 10^{-5}$, $U_3$ to $\sim 10^{-9}$, $U_4$ to $\sim 10^{-12}$, $U_6$ to machine precision; for the long $\rho_\Lambda$ descent ($n = 26$), $U_4$ is tight to $\sim 10^{-7}$ (\texttt{tools/research/lambda1\_power\_trace\_bounds.py}). The bound family thus determines $\lambda_1$ in closed form to any desired precision.
\end{corollary}

\begin{proof}
For positive eigenvalues $\lambda_1 \geq \lambda_2 \geq \cdots \geq 0$, $\mathrm{tr}\,C^k = \sum_j \lambda_j^k$, so $\lambda_1^k \leq \mathrm{tr}\,C^k$ and $\lambda_1 \leq U_k$. Monotonicity $U_k \geq U_{k+1}$ is the standard fact that $\ell^k$ norms of a positive vector decrease in $k$. The convergence $U_k \to \lambda_1$ at rate $(\lambda_2/\lambda_1)^k$ follows from $\mathrm{tr}\,C^k = \lambda_1^k(1 + (\lambda_2/\lambda_1)^k + \cdots)$ via the Bernoulli expansion of $(1 + \epsilon)^{1/k}$. The closed-form expression for $\mathrm{tr}\,C^k$ is the standard polynomial trace expansion; each $C_{ij}$ is a closed-form ratio of slicing ratios (Theorem~\ref{thm:gram-R} and its generalisation to arbitrary index pairs), so the sum is closed-form in cascade primitives.
\end{proof}

% [Empirical k_Q fitting section removed: superseded by the derived
% first-order exponential correction of Section~\ref{sec:first-order}
% below, whose numerical evidence is gathered in the table after
% Corollary~\ref{cor:exp-resum}. For three precision observables the
% Gram-matrix framework itself is superseded by the
% $\alpha(d^*)/\chi^k$ structural family of Paper~IVb, recorded in
% Remark~\ref{rem:superseded-by-family}.]

\subsection{Sign structure of the correction}

The geometric correction $\delta\Phi$ (Definition~\ref{def:delta-phi}) is strictly positive, since each adjacent overlap satisfies $C_{d,d+1}<1$ by Cauchy--Schwarz. Applied multiplicatively as $Q = Q_0\cdot\exp(\delta\Phi)$ (Corollary~\ref{cor:exp-resum}), it always increases the predicted value. This produces the observed sign structure across all cascade predictions:

\begin{itemize}[nosep]
\item \emph{Descent-dependent quantities} (those traversing multiple cascade layers) have negative leading-order deviations, because the independent-step approximation underestimates the propagated geometric content. The correction $\exp(+\delta\Phi)$ closes the gap.
\item \emph{Geometric quantities} (single-level ratios, topological constants) have positive leading-order deviations, because they do not traverse multiple layers and receive no inter-layer correction. Their deviations reflect the subleading structure of the Gamma function at a single level, which is a separate (and smaller) effect.
\end{itemize}

The clean sign separation---negative for descent, positive for geometric---is not an empirical observation. It is a consequence of $\delta\Phi > 0$, which is itself a consequence of $C_{ij} < 1$ for $i\neq j$ (Cauchy--Schwarz on non-collinear cascade integrands).

\subsection{The correction as compactification residual}

Theorem~4.2 establishes that each compactification step retains a residual $R_{\mathrm{eff}}(d) = 1/\sqrt{d+3}$. The overlap deficit (Theorem~\ref{thm:overlap}) scales as $1-C_{d,d+1}^2 \sim 1/(8d^2) \sim R_{\mathrm{eff}}^4/8$. The eigenvalue deficit $\epsilon$ is the cumulative effect of these per-step residuals over the path.

The independent-step approximation corresponds to the limit $R_{\mathrm{eff}} \to 0$ (complete compactification). The geometric correction $\delta\Phi$ (Definition~\ref{def:delta-phi}) is the leading-order cost of $R_{\mathrm{eff}} > 0$---the imperfect folding of the cascade's own geometry.

\subsection{The first-order correction: a Beta function identity}\label{sec:first-order}

\begin{theorem}[Adjacent overlap deficit correction]\label{thm:first-order}
For a descent-dependent observable traversing path $\{d_0,\ldots,d_0+n-1\}$:
\begin{equation}\label{eq:first-order}
\frac{\delta Q}{Q_0} = \sum_{k=0}^{n-2}\bigl(1 - C_{d_k,d_{k+1}}^2\bigr),
\end{equation}
where each $C_{d,d+1}$ is the adjacent-layer correlation (Theorem~\ref{thm:overlap}).
This is a sum of Beta function ratios. No free parameter enters.
\end{theorem}

\begin{proof}
The observable propagates through $n$ layers. The independent-step approximation treats each layer as perfectly collinear with the next ($C_{d,d+1}=1$). The first-order departure from collinearity at each step is the squared overlap deficit $1-C_{d,d+1}^2$, which measures the fraction of the integrand at step $d$ that is orthogonal to step $d+1$ (Theorem~\ref{thm:overlap}). For a multiplicative propagator, the leading correction is the sum of per-step deficits.
\end{proof}

\begin{corollary}[Exponential resummation forced by multiplicative composition]\label{cor:exp-resum}
For a multiplicative propagator $Q_0 = \prod_k (\text{step}_k)$
traversing the path of Theorem~\ref{thm:first-order}, the unique
closed-form leading-order Gram-corrected observable respecting
multiplicative composition under path splitting is
\begin{equation}\label{eq:exp-form}
Q \;=\; Q_0\cdot\exp\!\Bigl(\sum_{k=0}^{n-2}\bigl(1-C_{d_k,d_{k+1}}^2\bigr)\Bigr) \;+\; O(\delta^2),
\end{equation}
where $\delta = \max_k (1 - C_{d_k,d_{k+1}}^2)$. The exponential
matches Theorem~\ref{thm:first-order} at first order in $\delta$ per
step.
\end{corollary}

\begin{proof}
\emph{Setup.} Because the cascade recurrence is first-order
multiplicative (Corollary~3.2 of Paper~0), the Gram-corrected step at
layer $d_k$ receives a multiplicative factor
$1 + (1 - C_{d_k,d_{k+1}}^2) + O((1-C^2)^2)$, accumulating as a product:
\[
Q \;=\; Q_0\prod_{k=0}^{n-2}\bigl(1 + \delta_k + O(\delta_k^2)\bigr),
\qquad \delta_k := 1 - C_{d_k,d_{k+1}}^2.
\]
This product is the exact corrected observable (modulo $O(\delta^2)$
per-step residuals from Theorem~\ref{thm:first-order}'s first-order
truncation). Define the path-functional
$\Phi(P) := \sum_{k\in P} \delta_k$ and the correction factor
$K(P) := Q/Q_0$ over a path $P$.

\emph{Composition law.} Multiplicative composition under path
splitting is a structural property of the cascade recurrence: for a
path $P = P_1 \cup P_2$ formed by concatenating two consecutive
sub-paths, the corrected observable composes as
$Q(P) = Q_0(P_1)\cdot Q_0(P_2) \cdot K(P_1) \cdot K(P_2) = Q_0(P)\cdot K(P_1)\cdot K(P_2)$.
Therefore $K(P) = K(P_1)\cdot K(P_2)$, and since $\Phi(P) = \Phi(P_1) + \Phi(P_2)$
(the path-functional is additive under concatenation), $K$ as a
function of $\Phi$ must satisfy
\begin{equation}\label{eq:cauchy}
K(\Phi_1 + \Phi_2) \;=\; K(\Phi_1) \cdot K(\Phi_2).
\end{equation}

\emph{Cauchy's exponential functional equation forces the form.}
Equation~\eqref{eq:cauchy} together with continuity (a closed-form
correction is a smooth function of its path-functional argument) and
the boundary condition $K(0) = 1$ (no correction for an empty path)
has, by the standard solution of Cauchy's exponential functional
equation, the unique family of solutions $K(\Phi) = e^{c\Phi}$ for
some constant $c$. Matching the first-order behaviour to
Theorem~\ref{thm:first-order} ($K(\Phi) = 1 + \Phi + O(\Phi^2)$ for
small $\Phi$, hence $K'(0) = 1$) fixes $c = 1$. Therefore
\[
K(\Phi) \;=\; e^\Phi \;=\; \exp\!\Bigl(\sum_k \delta_k\Bigr),
\]
and the corrected observable is~\eqref{eq:exp-form}.

\emph{Why the linear form fails.} The linear closed-form
$K_{\mathrm{lin}}(\Phi) = 1 + \Phi$ matches at first order but
violates~\eqref{eq:cauchy}: $K_{\mathrm{lin}}(\Phi_1) K_{\mathrm{lin}}(\Phi_2)
= (1 + \Phi_1)(1 + \Phi_2) = 1 + \Phi_1 + \Phi_2 + \Phi_1\Phi_2 \neq
1 + (\Phi_1 + \Phi_2) = K_{\mathrm{lin}}(\Phi_1 + \Phi_2)$. The cross-term
$\Phi_1\Phi_2$ breaks additivity. The exponential is the unique
continuous closed form respecting composition; the linear form is the
unique closed form failing it.

\emph{Residual versus the exact product.} The exact corrected
observable is the product $Q_0\prod(1 + \delta_k + O(\delta_k^2))$, not
the exponential. The exponential and the product agree at first order
in $\delta$ per step but differ at second order: the exact product
expands as $1 + \sum\delta_k + \sum_{j<k}\delta_j\delta_k + O(\delta^3)$,
while the exponential expands as $1 + \sum\delta_k + \tfrac{1}{2}(\sum\delta_k)^2 + O(\delta^3)
= 1 + \sum\delta_k + \tfrac{1}{2}\sum\delta_k^2 + \sum_{j<k}\delta_j\delta_k + O(\delta^3)$.
The exponential overestimates the exact product by
$\tfrac{1}{2}\sum\delta_k^2$ at second order. The linear closed-form
$1 + \sum\delta_k$, by contrast, underestimates the product by
$\sum_{j<k}\delta_j\delta_k$ at second order.

For cascade paths starting at $d_0\geq 5$ the per-step deficits scale
as $\delta_k = 1 - C^2_{d_k,d_{k+1}} \sim 1/(8d^2)$
(Theorem~\ref{thm:overlap}), so
$\sum_{k}\delta_k^2 \lesssim \tfrac{1}{64}\sum_{d\geq 5} d^{-4} \approx
6\times 10^{-5}$, and the exponential overshoots the exact product by
at most $\tfrac{1}{2}(6\times 10^{-5}) \approx 3\times 10^{-5}$ in the
multiplicative correction (i.e.\ $\sim 0.003\%$ of $Q$). The linear
form underestimates the product by approximately
$\tfrac{1}{2}(\sum\delta)^2 \approx 2\times 10^{-4}$ ($\sim 0.022\%$
of $Q$) for the longest path ($d=5$ to $d=217$). The exponential is
therefore not only the unique composition-respecting closed form but
also the closer numerical approximation to the exact product than the
linear form, by an order of magnitude.

Both gaps are below the reported observational precisions of all
currently-tested cascade observables ($\sim 0.1\%$ or coarser). The
exponential is therefore the cascade's natural closed-form correction
--- forced by composition, exact at first order in $\delta$, and within
observational precision of the exact product for all paths considered.
It is the form used downstream (Paper~I, $\rho_\Lambda$ closure;
Paper~IVb, the $v$ and $m_\tau$ absolute-mass shifts).
\end{proof}

\begin{center}
\renewcommand{\arraystretch}{1.3}
\begin{tabular}{llcccccc}
\toprule
Observable & Path & $n$ & $\sum(1-C^2)$ & Leading & Corrected & Observed & Dev \\
\midrule
$\alpha_s$ & $d\!=\!5$--$12$ & 8 & 0.01186 & 0.1159 & 0.1173 & 0.1179 & $-0.5\%$ \\
$m_\tau/m_\mu$ & $d\!=\!6$--$13$ & 8 & 0.00938 & 16.53 & 16.69 & 16.82 & $-0.8\%$ \\
$m_\mu/m_e$ & $d\!=\!14$--$21$ & 8 & 0.00272 & 206.50 & 207.06 & 206.77 & $+0.1\%$ \\
$m_\tau$~(MeV) & $d\!=\!5$--$12$ & 8 & 0.01186 & 1755 & 1776 & 1777 & $-0.1\%$ \\
$v$~(GeV) & $d\!=\!5$--$12$ & 8 & 0.01186 & 240.8 & 243.7 & 246.2 & $-1.0\%$ \\
$\Omega_m^{\rm Bott}$ & $d\!=\!5$--$217$ & 213 & 0.02108 & 0.3115 & 0.3181 & 0.315 & $+1.0\%$ \\
\bottomrule
\end{tabular}
\end{center}

\noindent The first-order correction reduces descent deviations from $\sim\!2\%$ to sub-$1\%$ for all tested quantities. The correction for $m_\mu/m_e$ is small ($0.27\%$) because its path lies at higher $d$ where the per-step deficit scales as $1/(8d^2)$.

\begin{remark}[Scope of Theorem~\ref{thm:first-order}; status of $\Omega_m^{\rm Bott}$]\label{rem:sp9-status}
Theorem~\ref{thm:first-order} is proved for a \emph{multiplicative
propagator} traversing a contiguous path---i.e., for an observable of
the form $Q = \prod_{k=0}^{n-1}(\text{step factor at }d_k)$. Five of the
six table entries ($\alpha_s$, $m_\tau/m_\mu$, $m_\mu/m_e$, $m_\tau$,
$v$) have this structure directly. The sixth entry,
$\Omega_m^{\rm Bott}$, is structurally different: by
Theorem~5.10 of Paper~V (``Bott partition matter fraction''),
\[
\Omega_m^{\rm Bott}
\;=\;
\frac
{\displaystyle\sum_{\substack{d=5\\ d\bmod 8\in\{5,6\}}}^{217}\Omega_{d-1}}
{\displaystyle\sum_{d=5}^{217}\Omega_{d-1}},
\]
a \emph{ratio of sums} over cascade layers rather than a single
multiplicative propagator. Its inclusion in the table under the
full-path correction $\sum_{d=5}^{217}(1-C^2_{d,d+1}) = 0.02108$
applies the theorem outside its proven scope:
\begin{itemize}
\item The rigorous per-term treatment would correct each
$\Omega_{d-1}$ in the sums by its own path-length-dependent factor,
$\Omega_{d-1}\mapsto\Omega_{d-1}\cdot
\exp\bigl(\sum_{k=4}^{d-1}(1-C^2_{k,k+1})\bigr)$, then re-form the
ratio. The sums are dominated by the early terms near $d\sim 5$--$7$
(where $\Omega_d$ peaks), which have short paths and small corrections
($\sim\exp(0.001)\approx 1.001$). The resulting correction to the
ratio is therefore much smaller than the full-path factor
$\exp(0.02108)$.
\item The observed numerics support this: the leading value
$\Omega_m^{\rm Bott,\;lead}=0.31150$ sits $-1.1\%$ below the
observed $0.315$; the full-path ``corrected'' value $0.3181$
overshoots by $+1.0\%$. A rigorous per-term correction would fall
between these and is consistent with the two-population systematic
to leading order. The table entry reports the uniform full-path
correction for comparability with the other rows; it should be read
as an upper bound on the Gram shift for this observable, not as its
derived value.
\end{itemize}
A formal extension of Theorem~\ref{thm:first-order} to ratio-of-sums
observables---deriving the effective correction for
$Q=\sum_i a_i/\sum_j b_j$ where each $a_i,b_j$ is a multiplicative
propagator---is Paper~V Open Question~1 (and remains open: the
identification step from non-trivial-phase layers to matter is the
other half of that question). The $\Omega_m$ closure does not
\emph{require} the Gram correction: both the leading lapse identity
$\Omega_m=1/\pi=0.31831$ ($+1.1\%$, geometric population) and the
leading Bott value $0.31150$ ($-1.1\%$, descent-dependent population)
bracket observation to the cascade's stated precision.
\end{remark}

% [Derived coupling coefficient $k_Q^{(1)}$ and its observable-response
% remark removed: the empirical-fit linear parametrisation $Q = Q_0(1+k_Q\epsilon)$
% is superseded by the cascade-derived exponential form
% $Q = Q_0\cdot\exp(\sum(1-C^2))$ of Corollary~\ref{cor:exp-resum},
% which is uniquely forced by multiplicative composition under
% path-splitting. The auxiliary ratio $\sum(1-C^2)/\epsilon$ is no longer
% the structural object of interest; only the geometric shift
% $\delta\Phi=\sum(1-C^2)$ enters downstream usage.]

\subsection{What this section proves}

\begin{enumerate}[nosep]
\item The Gram matrix of cascade layer integrands is a matrix of Beta function values (Theorem~\ref{thm:gram}).
\item Adjacent layers have a strictly positive overlap deficit $1 - C_{d,d+1}^2 > 0$, scaling as $1/(8d^2)$ (Theorem~\ref{thm:overlap}). The correlation has the exact closed form $C^2_{d,d+1} = R(2d+2)^2 / [R(2d+1)\,R(2d+3)]$ in cascade slicing ratios (Theorem~\ref{thm:gram-R}), equivalently $\log C^2_{d,d+1} = -\tfrac{1}{2}\Delta^2 \log\alpha\big|_{2d+2}$ (Corollary~\ref{cor:gram-laplacian}): the Gram correction is the discrete log-Laplacian of the cascade scalar action's compliance function $\alpha(d) = R(d)^2/4$ at doubled argument, unifying the supplement with Paper~IVb's action principle (Remark~\ref{rem:gram-action-unification}).
\item The correlation matrix has one dominant eigenvalue $\lambda_1 \approx n(1 - \epsilon)$ with $\epsilon > 0$ (Theorem~\ref{thm:eigenvalue}).
\item The eigenvalue deficit is derived from first-order perturbation theory as a sum of Beta function deficits (Theorem~\ref{thm:pert}).
\item The sign of the correction (positive) explains the clean sign separation between descent and geometric predictions.
\item The first-order multiplicative correction $\delta Q / Q_0 = \sum(1-C^2_{d,d+1})$ (Theorem~\ref{thm:first-order}), uniquely promoted to the exponential form $Q = Q_0 \cdot \exp(\sum(1-C^2))$ by multiplicative composition under path-splitting (Corollary~\ref{cor:exp-resum}). Both ingredients are derived; no free parameter enters.
\item Every quantity in this section is a Gamma function identity. No physics enters.
\end{enumerate}

\subsection{Open questions}\label{sec:open}

\begin{enumerate}[nosep]
\item \emph{Analytic formula for $\epsilon(n, d_0)$.} Partially closed. At $n = 2$ the dominant eigenvalue admits the exact closed form $\lambda_1 = 1 + R(2d_0+2)/\sqrt{R(2d_0+1)R(2d_0+3)}$ (Corollary~\ref{cor:lambda1-n2}). For $n \geq 3$ the power-trace bound family $U_k = (\mathrm{tr}\,C^k)^{1/k}$ (Corollary~\ref{cor:lambda1-power-trace}) provides closed-form cascade-native approximations to $\lambda_1$ converging at geometric rate $(\lambda_2/\lambda_1)^k$, reaching machine precision at $k = 6$ for cascade-physics paths. What remains genuinely open is whether a single-equation radical or special-function closed form (rather than a convergent family of upper bounds) exists at $n \geq 3$.
\item \emph{Higher-order corrections.} The subdominant eigenvalues $\lambda_3, \lambda_4, \ldots$ contribute at third order and beyond. For the $d = 5, \ldots, 12$ path, $\lambda_2/n = 0.81\%$ and $\lambda_3/n = 0.008\%$: the third-order correction is negligible. For longer paths, it may not be.
\item \emph{Application to mixed quantities.} The first-order correction (Theorem~\ref{thm:first-order} and Corollary~\ref{cor:exp-resum}) applies cleanly to quantities with a single unambiguous multiplicative cascade path. Quantities expressible as ratios of sums over cascade layers ($\Omega_m^{\rm Bott}$ in Paper~V) require a per-term treatment articulated in Remark~\ref{rem:sp9-status} but not formally derived. A formal extension of the first-order theorem to such ratio-of-sums observables is open.
\end{enumerate}

\begin{remark}[Status of the correction]
The first-order correction is derived from cascade-native first principles. Theorem~\ref{thm:gram-R} and Corollary~\ref{cor:gram-laplacian} establish that the per-step Gram correction is the discrete log-Laplacian of the cascade scalar action's compliance function $\alpha(d)=R(d)^2/4$, evaluated at the doubled cascade argument. Theorem~\ref{thm:first-order} sums this to the path-cumulative shift $\delta\Phi=\sum(1-C^2_{d,d+1})$, and Corollary~\ref{cor:exp-resum} promotes it to the unique closed form $Q = Q_0\cdot\exp(\delta\Phi)$ respecting multiplicative composition under path-splitting. Both single-source $\alpha(d^*)/\chi^k$ family of Paper~IVb and the path-distributed Gram correction here therefore derive from the same compliance $\alpha(d)=R(d)^2/4$: source response and discrete log-Laplacian, respectively, of one structural object.
\end{remark}

\begin{remark}[Superseded for the eight precision observables]\label{rem:superseded-by-family}
For eight Standard Model precision observables---$\alpha_s(M_Z)$, $m_\tau/m_\mu$, $m_\tau$ absolute, $\ell_A$, $\sin^2\theta_W$, $\Omega_m$, $\theta_C$, and $b/s$---the Gram-matrix framework above is superseded by the sharper structural identity of Paper~IVb. Both deviations in these observables are closed \emph{at experimental precision} by cascade potential shifts of the form $\delta\Phi = \alpha(d^*)/\chi^k$, where $d^*$ is a distinguished dimension from this paper's four-dimension tower (Theorem~7.1) and $\chi=\chi(S^{2n})=2$ is the Euler characteristic:
\begin{itemize}
\item $\alpha(14)/\chi$ closes $\alpha_s$ and $m_\tau/m_\mu$ together (Part~IVb, Theorem~4.3);
\item $\alpha(19)/\chi$ closes the $\tau$ absolute mass and $\ell_A$ at the phase-transition threshold $d_1=19$ (Part~IVb, Theorem~4.4);
\item $\alpha(5)/\chi^3$ closes $\sin^2\theta_W$ and $\Omega_m$ at the volume maximum $d_V=5$;
\item $\alpha(7)/\chi^2$ closes $\theta_C$ and $\alpha(7)/\chi^4$ closes $b/s$, both sourced at the area maximum $d_0=7$.
\end{itemize}
For these eight observables the cascade closes at experimental precision via a single-source mechanism that does not invoke the Gram framework at all. The first-order Gram correction here remains the appropriate cascade-native account for the remaining descent-dependent observables ($m_\mu/m_e$ to $0.13\%$, $v$ before its non-perturbative $\exp(-\pi/\alpha(5))$ closure, etc.), where it reduces deviations from $\sim 2\%$ to sub-$1\%$.
\end{remark}

\end{document}
